% TOP_PROBLEM: {"id": 30006680, "title": "Tate Conjecture for Algebraic Cycles", "category_id": 5, "category": {"id": 5, "name": "algebraic_geometry", "display_name": "Algebraic Geometry", "description": "Geometric objects defined by polynomial equations.", "slug": "algebraic-geometry", "order_index": 5, "created_at": "2026-07-31T15:26:25.671Z"}, "status": "open"}
% FIRST_POSTED: null
% !TeX program = lualatex
% Compile twice: lualatex open_problems_500.tex
% All references and prose are embedded; no BibTeX database or external inputs.
\documentclass[11pt,a4paper]{article}
\usepackage[margin=24mm,headheight=14pt]{geometry}
\usepackage{fontspec}
\setmainfont{Latin Modern Roman}
\setsansfont{Latin Modern Sans}
\setmonofont{Latin Modern Mono}
\usepackage{amsmath,amssymb,mathtools}
\usepackage{unicode-math}
\setmathfont{Latin Modern Math}
\usepackage{microtype}
\usepackage{enumitem,xurl,needspace}
\usepackage[dvipsnames]{xcolor}
\usepackage{hyperref}
\hypersetup{unicode=true,colorlinks=true,linkcolor=MidnightBlue,urlcolor=MidnightBlue,
 pdftitle={Mathematical Research Notebook},pdfauthor={Source-annotated compilation},bookmarksnumbered=true}
\usepackage{bookmark}
\usepackage{tocloft}
\setlength{\cftsecnumwidth}{3em}
\cftsetpnumwidth{2.5em}
\cftsetrmarg{4em}
\usepackage{titlesec}
\titleformat{\section}{\Large\bfseries\raggedright}{\thesection}{0.7em}{}
\usepackage{fancyhdr}
\pagestyle{fancy}\fancyhf{}
\fancyhead[L]{\small\sffamily Open Problems Math | Research notebook}
\fancyhead[R]{\small 27 September 2026}
\fancyfoot[C]{\thepage}
\setlength{\parindent}{0pt}
\setlength{\parskip}{4pt plus 1pt}
\setlength{\emergencystretch}{3em}
\setcounter{tocdepth}{1}
\setcounter{secnumdepth}{1}
\setlist[itemize]{leftmargin=3em,itemsep=3pt,topsep=4pt,parsep=0pt}
\allowdisplaybreaks
\newcommand{\N}{\mathbb{N}}
\newcommand{\Z}{\mathbb{Z}}
\newcommand{\Q}{\mathbb{Q}}
\newcommand{\R}{\mathbb{R}}
\newcommand{\C}{\mathbb{C}}
\newcommand{\F}{\mathbb{F}}
\newcommand{\A}{\mathbb{A}}
\newcommand{\PP}{\mathbb P}
\newcommand{\E}{\mathbb{E}}
\newcommand{\eps}{\varepsilon}
\newcommand{\dd}{\,\mathrm d}
\newcommand{\sref}[2]{\hyperlink{source:#1:#2}{[S#2]}}
\newcommand{\eref}[2]{\hyperlink{extra:#1:#2}{[E#2]}}
% Turn the next switch off for a shorter reading copy. The quotations preserve
% every exactTarget field, including qualifications omitted from a short summary.
\newif\ifshowcatalogue
\showcataloguetrue
\newcommand{\cataloguescope}[1]{\ifshowcatalogue
 \paragraph{Exact scope in the supplied catalogue.}
 \begin{quote}\small #1\end{quote}\fi}
\usepackage{etoolbox}
\AtBeginEnvironment{itemize}{\small}
% Standalone entry; original section content is delimited for verification.
\begin{document}
\setcounter{section}{0}
%% BEGIN ORIGINAL SECTION CONTAINER
% ================================================================
% problemId: 30006680
% Canonical problem ID: 30006680
% exactTarget SHA-256: 8b407cd86d19b5329407c4da4c848f802712f0ba624bf22a9861a26898da7064
\clearpage
\section*{\texorpdfstring{Tate Conjecture for Algebraic Cycles}{Tate Conjecture for Algebraic Cycles}}\label{30006680}
\subsection{Definitions and mathematical statement}
Let $\bar X=X\times_k\bar k$, $G_k=\operatorname{Gal}(\bar k/k)$, and $\ell\ne\operatorname{char}k$. The Chow group $\operatorname{CH}^i(X)$ consists of codimension-$i$ cycles modulo rational equivalence. The Tate twist $\Q_\ell(i)$ is the $i$-fold tensor power of the cyclotomic representation. The assertion is surjectivity of
\[
 \operatorname{cl}_\ell:\operatorname{CH}^i(X)\otimes\Q_\ell
 \longrightarrow H^{2i}_{\mathrm{\acute et}}(\bar X,\Q_\ell(i))^{G_k}
\]
for every smooth geometrically irreducible projective $X/k$, every $i$, and every $k$ finitely generated over its prime field. The domain uses cycles over $k$, and the coefficients are $\Q_\ell$, not $\Z_\ell$.
\par\smallskip\textit{Formulation and concept references:} \sref{30006680}{1}.
\cataloguescope{For every smooth geometrically irreducible projective variety X over a field k finitely generated over its prime field, every codimension i, and every prime l different from char(k), prove that the l-adic cycle-class map from codimension-i algebraic cycle classes defined over k, tensored with Q\_l, is surjective onto H\textasciicircum{}(2i)(X\_bar,Q\_l(i))\textasciicircum{}Gal(k\_bar/k).}
\subsection{Short English statement}
Are all arithmetic-symmetry-invariant even cohomology classes generated by algebraic cycles defined over the base field?
%% BEGIN RESEARCH INSERTION
\subsection{Research attempt}
\paragraph{Current frontier.}
\textit{Research check: 27 September 2026.}
The divisor theorem for abelian varieties is an established special case;
it does not supply all higher-codimension classes on arbitrary varieties.
The distinction between cycle-class surjectivity and the stronger package
including semisimplicity is explicit in Milne, Section 1, especially
Theorem 1.4. \sref{30006680}{1}, \eref{30006680}{1}.
Kahn's 2025 surface manuscript gives a reformulation and identifies a step
that does not yet work in Section 3.2. Milne's September 2025 note concerns
failure with \emph{integral} coefficients, not a counterexample to the
rational statement here. \sref{30006680}{4}, \sref{30006680}{3}.
We investigate a finite-dimensional certificate and the descent needed to
make its cycles live over the stipulated base field.

\paragraph{Attack route.}
One route is to lift invariant classes from special fibers; it requires
algebraic cycles to lift, not merely their cohomology. Another is to build
cycles until their intersection matrix exhausts the Frobenius eigenspace.
A third is to find cycles after a field extension and descend them. The
second and third routes can be combined without assuming the conjectural
semisimplicity one is trying to exploit. The finite-field criterion below
is stronger than the selected surjectivity statement; this strength is
kept explicit.

\paragraph{Attempt: descent of an already algebraic invariant class.}
Let $L/k$ be finite Galois, $p:X_L\to X$ the projection, and
$\alpha\in H^{2i}(\bar X,\Q_\ell(i))^{G_k}$. Suppose that
$\alpha=\operatorname{cl}_\ell(z)$ for
$z\in\operatorname{CH}^i(X_L)\otimes\Q_\ell$.
Then the cycle
\[
 z_0=\frac{1}{[L:k]}p_*z\in\operatorname{CH}^i(X)\otimes\Q_\ell
 \tag{15.1}
\]
has class $\alpha$. Indeed, on geometric cohomology, pullback of the
pushforward is $\sum_{\sigma\in\operatorname{Gal}(L/k)}\sigma$; every
summand fixes $\alpha$. Geometric base change identifies the target
cohomology groups, so division by $[L:k]$ proves the claim. This uses
rational $\ell$-adic coefficients even when $\ell$ divides the degree.
It is not an integral argument. The push-pull and cycle-class
compatibilities used here are those of the formulation and of the
finite-field discussion in \eref{30006680}{1}.

Consequently, if a $G_k$-invariant class is represented by a finite
combination of cycles defined over a finite separable extension, its
\emph{field of definition} is not an additional obstruction: take a
Galois closure and apply (15.1). This does not construct the cycles.
In positive characteristic, this particular lemma is stated only for
separable extensions; no unmentioned passage through a purely inseparable
extension is needed for it.

\paragraph{Attempt: a certificate that also controls Jordan blocks.}
Now specialize to $k=\F_q$ and write $n=\dim X$,
\[
 V=H^{2i}(\bar X,\Q_\ell(i)),\qquad
 W=H^{2n-2i}(\bar X,\Q_\ell(n-i)).
\]
Let $F$ be geometric Frobenius acting on these \emph{twisted} spaces.
Poincare duality gives a nondegenerate $F$-invariant pairing
$\langle\ ,\ \rangle:V\times W\to\Q_\ell$.
Let $r$ be the algebraic multiplicity of $1$ in the characteristic
polynomial of $F$ on $V$. Assume we have cycles over $k$,
\[
 z_1,\ldots,z_r\in\operatorname{CH}^i(X)_{\Q},\qquad
 w_1,\ldots,w_r\in\operatorname{CH}^{n-i}(X)_{\Q},
\]
such that the rational matrix
\[
 D_{ba}=\deg(z_a\cdot w_b)\quad(1\leq a,b\leq r)
 \quad\hbox{satisfies}\quad\det D\ne0. \tag{15.2}
\]
Then the selected Tate assertion holds in codimension $i$ for $X$,
and $F$ acts as the identity on its entire generalized $1$-eigenspace.

Here is a proof, including the dimension issue. The classes
$v_a=\operatorname{cl}(z_a)$ lie in $V^F$. Pairing a relation
$\sum_a t_av_a=0$ with all $\operatorname{cl}(w_b)$ yields $Dt=0$,
so they are independent. On the other hand
$\dim V^F\leq r$, because the fixed space is contained in the generalized
$1$-eigenspace. Thus $\dim V^F=r$ and these classes span it.
Equality with the generalized eigenspace excludes all nontrivial Jordan
blocks at $1$. No semisimplicity was used as an input. If $r=0$,
the fixed space is zero and the assertion is vacuous.

More explicitly, the algebraic correspondence
\[
 \Pi=\sum_{a,b}(D^{-1})_{ab}\,w_b\times z_a
       \in\operatorname{CH}^n(X\times X)_{\Q} \tag{15.3}
\]
acts on $V$ by
$\Pi_*(v)=\sum_{a,b}v_a(D^{-1})_{ab}\langle v,\operatorname{cl}(w_b)\rangle$.
It is the identity on the span of the $v_a$, has image that span, and
therefore satisfies $\Pi_*^2=\Pi_*$ on $V$. In fact this particular
product-cycle formula gives an idempotent even modulo rational
equivalence. The composition of $w_d\times z_c$ followed by
$w_b\times z_a$ is $D_{bc}w_d\times z_a$: in the middle factor the
intersection is a zero-cycle whose pushforward is its degree. Summing
with the coefficients in (15.3) gives the matrix product
$D^{-1}DD^{-1}=D^{-1}$, hence $\Pi\circ\Pi=\Pi$ in
$\operatorname{CH}^n(X\times X)_{\Q}$. No assertion that arbitrary
cohomological projectors lift to Chow projectors is used; this one is
built from the assumed dual cycle lists. The connection with the
finite-field Tate and semisimplicity conditions is discussed in
\eref{30006680}{1}, Theorem 1.4; no historical novelty is asserted.

\paragraph{Attempt: what approximation would actually suffice.}
Let $A\subseteq V$ be the $\Q_\ell$-span of all cycle classes over $k$.
Because $V$ is finite dimensional, $A$ is a closed subspace. Thus an
actual sequence $a_j\in A$ converging $\ell$-adically to $\alpha$ would
establish $\alpha\in A$. To see closedness directly, choose a basis of
$A$, extend it to one of $V$, and write $A$ as the common kernel of the
remaining continuous coordinate functionals. The subtlety is not this
limit; it is producing the approximants from actual cycles, with
compatible accuracy and the correct base field.

\paragraph{Stress test.}
Two elementary models obstruct tempting replacements for that missing
construction. On the lattice $\Z_\ell$, multiplication by $1+\ell$ fixes
every class modulo $\ell$, but its fixed space on $\Q_\ell$ is zero.
Mod-$\ell$ invariance is therefore not a certificate of a rational Tate
class. Independently, take
\[
 F=\begin{pmatrix}1&1\\0&1\end{pmatrix},\qquad
 F^\vee=(F^{-1})^{\mathsf T}.
\]
The fixed line of $F$ is spanned by $e_1$; the fixed line of $F^\vee$ is
spanned by the second coordinate functional. Their pairing is zero.
All eigenvalues equal $1$, yet the invariant pairing is degenerate and
(15.2) with $r=2$ is impossible. This refutes the purely linear-algebraic
inference from eigenvalues of the correct weight to semisimplicity. It
is not asserted to be the cohomology of a counterexample variety.

If cycles are first constructed over an extension, their individual
classes need not be fixed by $G_k$; averaging can kill noninvariant
components and lower rank. Formula (15.1) avoids that loss only for the
specified invariant class already represented there. Finally, neither
point counts determining $r$ nor a positive numerical intersection form
produces $r$ independent algebraic classes. For finitely generated fields
other than finite fields, the one-operator certificate does not replace
the full $G_k$-invariance condition. These are genuine limits of the route.

\paragraph{Outcome.}
\textbf{Outcome: Conditional reduction.}
A nondegenerate cycle-intersection matrix of the full Frobenius
multiplicity would prove the finite-field case under consideration and
supply an explicit algebraic idempotent correspondence. An invariant class already
algebraic over a finite separable extension descends with rational
coefficients. Neither step establishes the required abundance of cycles,
and the matrix criterion deliberately includes an additional
semisimplicity conclusion. It is not claimed equivalent to, or weaker
than, the selected conjecture. Historical novelty of these auxiliary
observations is not asserted.
%% END RESEARCH INSERTION
\subsection{Sources}
\begin{itemize}\raggedright
\item[\textnormal{[S1]}]\hypertarget{source:30006680:1}{} Li, What is the Tate conjecture?, map (2) and Conjecture 1 (2013).
\url{https://www.math.columbia.edu/~chaoli/docs/TateConjecture.html}.
{\small\textit{Catalogue formulation source. Consulted 24 September 2026: Cycle-class map and Conjecture 1.}}
\item[\textnormal{[S2]}]\hypertarget{source:30006680:2}{} An approach to the Tate conjecture for surfaces over a finite field — Bruno Kahn (2025).
\url{https://arxiv.org/abs/2505.07337}.
{\small\textit{Catalogue research/status reference; not a proof endorsement.}}
\item[\textnormal{[S3]}]\hypertarget{source:30006680:3}{} On the integral Tate conjecture for abelian varieties — J. S. Milne (September 8, 2025).
\url{https://arxiv.org/abs/2509.06707}.
{\small\textit{Catalogue research/status reference; not a proof endorsement.}}
\item[\textnormal{[S4]}]\hypertarget{source:30006680:4}{} An approach to the Tate conjecture for surfaces over a finite field — Bruno Kahn, full text.
\url{https://arxiv.org/html/2505.07337v1}.
{\small\textit{Catalogue research/status reference; not a proof endorsement.}}
%% BEGIN RESEARCH REFERENCES
\item[\textnormal{[E1]}]\hypertarget{extra:30006680:1}{} James S. Milne, \emph{The Tate conjecture over finite fields}, version 2.1, 7 November 2007, Section 1 and Theorem 1.4.
\url{https://www.jmilne.org/math/articles/2007e.pdf}.
{\small\textit{Added research source. Cycle surjectivity, intersection pairing and the stronger semisimplicity package. Consulted 27 September 2026.}}
%% END RESEARCH REFERENCES
\end{itemize}

%% END ORIGINAL SECTION CONTAINER
\end{document}
